\section{Основная часть}
\label{sec:main}

Нейроны кодируют информацию, избирательно представляя различные стимулы, поведенческие паттерны или когнитивные состояния посредством своей активности, формируя тем самым основу нейронного кода \cite{DayanAbbott2001}. На популяционном уровне эти сигналы представляют собой многомерные переменные, определяющие поведение и когнитивные процессы \cite{Averbeck2006}. Вместе с тем при работе со свободно ведущими себя животными центральная трудность состоит не столько в обнаружении связей между нейронной активностью и поведением как таковых, сколько в корректной атрибуции активности нейрона конкретной объясняющей переменной в условиях, когда поведение является многомерным, характеризуется взаимными корреляциями и обладает выраженной временн\'{о}й структурой.

Нейронная селективность динамически формируется под влиянием опыта. В гиппокампе клетки места способны быстро образовывать пространственные рецептивные поля \cite{Sotskov2022}, причём в ряде случаев эти карты остаются стабильными на протяжении недель \cite{Wilson1993, Thompson1990}. В то же время широко распространена смешанная селективность \cite{TyeMiller2024}: один нейрон может отражать комбинации нескольких стимулов, поведенческих или когнитивных переменных, а не какой-либо один признак. Так, нейроны префронтальной коры формируют смешанную селективность к задачно-релевантным переменным в ходе обучения \cite{mix3, Mante2013}. Даже канонические примеры функционально специализированных нейронов могут кодировать множественные переменные: клетки места зачастую отражают не только местоположение, но и время, вознаграждение и траектории \cite{Zheng2021, Lian2022}. Этот принцип распространяется и на сенсорную обработку: нейроны первичной зрительной коры и ретросплениальной коры способны кодировать как запомненные, так и воспринимаемые в данный момент изображения \cite{Makino2015}. Совокупность этих данных указывает на то, что наблюдаемую селективность необходимо интерпретировать с учётом других коррелированных переменных, совпадающих по времени с исследуемой в условиях естественного поведения.

Проблема атрибуции становится особенно острой в натуралистических условиях, поскольку поведенческие переменные по своей природе коррелированы. Развитие методов автоматического анализа поведения --- как неконтролируемых (MoSeq) \cite{Weinreb2024}, так и контролируемых (DeepAction \cite{Harris2023deepaction}, MARS \cite{Segalin2021}) и основанных на правилах (BehaviorDEPOT) \cite{Gabriel2022} --- позволяет извлекать десятки и сотни поведенческих признаков из видеозаписей, порождая богатые, но сильно коррелированные наборы данных. Это создаёт проблему идентифицируемости: кажущаяся селективность к одной переменной может объясняться селективностью к коррелированному ковариату. Например, активность нейрона может выглядеть как кодирование локомоции, тогда как на самом деле настройка обусловлена скоростью; аналогично, ответ клетки места может быть ошибочно интерпретирован как реакция на социальное взаимодействие лишь потому, что встречи приурочены к определённым местам. Таким образом, широкая распространённость смешанной селективности порождает фундаментальную проблему интерпретации. Эта трудность характерна для нейронаук в целом: ненаправленные движения объясняют значительную долю активности по всей коре \cite{Musall2019}, тета-модуляция затрудняет разделение кодирования скорости и фазы в гиппокампе \cite{Vanderwolf1969, McNaughton1983}. Приведённые примеры подчёркивают общее требование к анализу селективности: необходимо явно учитывать коррелированные ковариаты и временн\'{у}ю структуру нейронных и поведенческих сигналов.

Кальциевый имиджинг дополнительно усложняет эту задачу, обеспечивая крупномасштабную регистрацию от сотен до тысяч нейронов в ходе непрерывного поведения \cite{Chen2013, Zong2022}. В отличие от спайков, кальциевый флуоресцентный сигнал является непрерывным и характеризуется сложной временн\'{о}й динамикой. Даже самые быстрые индикаторы (например, jGCaMP8) демонстрируют времена затухания от десятков до сотен миллисекунд \cite{Zhang2023}. Кинетика индикатора приводит к временн\'{о}й свёртке, при которой флуоресцентный сигнал от более ранних событий ложно коррелирует с последующим поведением \cite{Wei2020}. Кроме того, поведение описывается переменными принципиально различных типов --- от дискретных категориальных состояний (груминг, стойка на задних лапах) до непрерывных величин (скорость, направление головы), --- что требует аналитических подходов, способных работать с гетерогенными типами переменных в рамках единой системы. Многие методы анализа спайков предполагают точечно-процессные наблюдения и не переносятся непосредственно на темпорально фильтрованные флуоресцентные сигналы; в то же время деконволюция может вносить модельно-зависимые искажения и снижать качество последующего декодирования или атрибуционного анализа \cite{Rupprecht2021}.

Многие из существующих подходов удовлетворяют лишь части перечисленных требований. Линейные меры не способны обнаруживать нелинейные зависимости, снижение размерности скрывает вклад отдельных нейронов, а вероятностные декодеры могут предсказывать сложное поведение, не позволяя при этом количественно определить, какие именно поведенческие переменные объясняют активность данного нейрона, --- особенно в условиях коррелированных ковариатов \cite{Kriegeskorte2019}.
Информационно-теоретические подходы предоставляют принципиальный формализм для количественной оценки связей между нейронной активностью и поведением на непрерывных данных, однако существующие программные инструменты ориентированы преимущественно на иные экспериментальные режимы или задачи анализа. Наиболее полный из них --- Neuroscience Information Toolbox (NIT) \cite{Maffulli2022} --- предлагает множество оценок взаимной информации с коррекцией смещения и подробными рекомендациями для кальциевого имиджинга, однако рассчитан на парадигмы с повторными пробами (trial-based), предполагающие большое число повторений для каждого условия. Напротив, непрерывные записи свободного поведения требуют сохранения временн\'{о}й автокорреляции, оценки оптимальных задержек между активностью нейрона и поведением, а также контроля ковариации поведенческих признаков. Другие программные пакеты ориентированы на дополняющие уровни анализа: MINT --- на передачу информации на популяционном уровне \cite{MINT2025}; FRITES --- на групповой вывод в электрофизиологических данных \cite{Frites2022}; IDTxl --- на анализ сетевой связности \cite{IDTxl2019}; HOI --- на анализ взаимодействий высших порядков \cite{Neri2024}. Метод оценки взаимной информации на основе гауссовой копулы (Gaussian Copula Mutual Information, GCMI) \cite{gcmi} обеспечивает эффективное вычисление взаимной информации для непрерывных данных, однако сама по себе оценка взаимной информации не решает ключевых проектировочных задач, возникающих при свободном поведении: выбора задержки в условиях автокорреляции и разделения истинного кодирования и артефактов, обусловленных корреляциями. Таким образом, сохраняется пробел: ни один из существующих инструментов не обеспечивает информационно-теоретический анализ селективности, который одновременно работал бы с непрерывным кальциевым флуоресцентным сигналом, переменными смешанных типов, оптимизацией временных задержек и разделением истинного кодирования и корреляционных артефактов --- в масштабе тысяч нейронов и десятков поведенческих признаков.

В настоящей работе представлен INTENSE (INformation-Theoretic Evaluation of Neuronal SElectivity) --- инструмент с открытым исходным кодом, объединяющий эффективную оценку взаимной информации с компонентами анализа, адаптированными для непрерывного кальциевого имиджинга при свободном поведении. INTENSE обеспечивает информационно-теоретическую оценку селективности отдельных нейронов в большом масштабе с учётом временных задержек, гетерогенных поведенческих переменных и коррелированного пространства поведенческих признаков.

